\chapter{Supplementary Evaluation Material}

This appendix collects the concrete sample, prompt, configuration, and launch details that support the main results chapters. The goal is traceability: when the main text refers to a sample family, a variant label, or a maintained preset, the corresponding evaluator-side details should be easy to find here without interrupting the flow of the main argument.

\section{Sample and Variant Traceability}
\label{app:sample-variant-traceability}

The thesis uses three evaluation groups with different roles. The broad pilot screens a wider medium-and-hard authored corpus for promising workflow dimensions. The depth-focused studies reuse a smaller decoder-oriented slice to stress stripped and packed recovery tasks. The final held-out evaluation uses five opaque aliases backed by evaluator-side references and is the strongest basis for the main empirical claims. Tables~\ref{tab:broad-sweep-samples}--\ref{tab:final-heldout-samples} map each executable group to the implemented methods or intended analysis challenge it represents, and Table~\ref{tab:evaluation-variants} maps the maintained variant labels back to their baseline-relative changes.

\subsection{Broad Coverage Pilot Executables}

The broad pilot preset \texttt{coverage\_broad\_r1\_60usd} runs over the medium-and-hard authored slice of the experimental corpus. Its purpose is not to emulate a single malware family but to cover a range of reverse-engineering challenges including decoy-heavy naming, stripping, hashing-based dispatch, packing, embedded payload parsing, and anti-analysis handling.

\small
\begin{longtable}{p{0.26\textwidth} p{0.09\textwidth} p{0.57\textwidth}}
\caption{Executable set used by the maintained broad-coverage pilot. The third column summarizes the main implemented techniques or intended analysis challenge for each sample.}
\label{tab:broad-sweep-samples}\\
\textbf{Sample} & \textbf{Diff.} & \textbf{Major implemented methods or intended challenge} \\
\hline
\endfirsthead
\textbf{Sample} & \textbf{Diff.} & \textbf{Major implemented methods or intended challenge} \\
\hline
\endhead
\texttt{maintenance\_orchestrator\_test.exe} & hard & Misleading function names, decoy strings, dormant decoded notices, and a live/dead opcode dispatcher wrapped around a deterministic local maintenance-report workflow; primarily tests false-positive control and live-path recovery. \\
\texttt{maintenance\_orchestrator\_test\_stripped.exe} & hard & Stripped version of the maintenance-orchestrator case; removes symbol-name shortcuts and forces the same benign workflow to be recovered from code and data flow alone. \\
\texttt{callback\_dispatch\_test.exe} & medium & Function-pointer callback table plus struct-based command descriptors and indirect calls; tests recovery of a callback-driven state machine from data structure and dispatch evidence. \\
\texttt{config\_decoder\_test.exe} & medium & Single-byte XOR configuration blob, key/value parser, stack-buffer decode path, and validation branching; serves as the clean decoder baseline for later depth-focused studies. \\
\texttt{signal\_router\_report\_test.exe} & medium & Threat-looking helper names and decoy labels hide local signal-routing and report-generation behavior; tests whether the workflow follows real data flow instead of surface wording. \\
\texttt{stack\_notice\_scheduler\_test.exe} & medium & Stack-built decoy strings, dormant XOR-decoded notices, decoded minute offsets, and insertion-sort schedule ordering; tests cautious interpretation of FLOSS-style recovered strings and scheduling logic. \\
\texttt{multilayer\_encode\_test.exe} & hard & Layered ROL+XOR encoding around URLs, commands, per-entry metadata, and a magic header; tests multi-stage decode reasoning rather than one-pass string recovery. \\
\texttt{hash\_dispatch\_test.exe} & hard & DJB2 hash algorithm, hash-to-handler dispatch table, stack-built command strings, and dead hash entries; tests concrete hash-dispatch reconstruction. \\
\texttt{hash\_dispatch\_test\_stripped.exe} & hard & Stripped variant of the hash-dispatch case; preserves the dispatcher logic while removing preserved names, so behavior must be recovered from constants and control flow. \\
\texttt{embedded\_payload\_test.exe} & hard & Embedded blob with a structured header, rolling checksum, per-record XOR keys, and staged extraction logic; tests payload parsing and staged extraction reasoning. \\
\texttt{embedded\_payload\_test\_upx.exe} & hard & UPX-packed version of the embedded-payload case; adds packer recognition and unpack-aware reasoning before the same staged extraction problem. \\
\texttt{branch\_weave\_snapshot\_test.exe} & hard & Encoded operation stream and dispatcher hidden behind misleading names, weak live strings, and inert decoded notices; tests control-flow-led behavior recovery when strings are poor evidence. \\
\texttt{branch\_weave\_snapshot\_test\_stripped.exe} & hard & Stripped version of the branch-weave case; pushes the same grid-building and report-generation logic into a symbol-poor setting. \\
\texttt{anti\_analysis\_suite\_test.exe} & hard & Anti-debug APIs, timing probes, environment fingerprinting, rotating-key XOR config, dead code, and misleading names; tests anti-analysis identification without inventing unsupported payload claims. \\
\hline
\end{longtable}
\normalsize

\subsection{Focused Decoder-Depth Executables}

The repeated decoder-depth preset \texttt{sweep\_decoder\_depth\_r3} and the follow-up preset \texttt{sweep\_decoder\_depth\_followups\_r3} both reuse the same two executables and the same deeper task slate: \texttt{default\_analysis}, \texttt{config\_value\_recovery}, \texttt{decode\_parser\_flow\_recovery}, \texttt{parser\_validation\_audit}, and \texttt{packing\_decoder\_triage}. This narrower group is intended to stress symbol stripping, parser reconstruction, and packer-aware recovery more heavily than the broad pilot does.

\begin{table}[p]
\centering
\small
\begin{tabular}{p{0.30\textwidth} p{0.10\textwidth} p{0.50\textwidth}}
\textbf{Sample} & \textbf{Diff.} & \textbf{Major implemented methods or intended challenge} \\
\hline
\texttt{config\_decoder\_test\_stripped.exe} & hard & Stripped decoder sample with single-byte XOR configuration recovery, parser logic, and validation branching; stresses recovery of decode and parser behavior without preserved symbol names. \\
\texttt{config\_decoder\_test\_upx\_stripped.exe} & hard & Same underlying decoder logic under UPX packing and stripped symbols; adds packer triage and unpack-aware reasoning before control-flow and data-flow recovery. \\
\end{tabular}
\caption{Executable set used by the repeated and follow-up decoder-depth studies.}
\label{tab:depth-sweep-samples}
\end{table}

\subsection{Final Held-Out Executables}

The final preset \texttt{thesis\_final\_eval\_r3\_100usd} runs over five opaque aliases under \path{Testing/sources/final\_round/}. The runtime only sees the aliases in the first column. The evaluator-side family labels and behavior summaries remain reference-only metadata and are included here so the results chapter can discuss held-out sample heterogeneity precisely without leaking those labels into the prompt surface itself.

\begin{table}[p]
\centering
\small
\begin{tabular}{p{0.14\textwidth} p{0.19\textwidth} p{0.53\textwidth}}
\textbf{Sample} & \textbf{Evaluator family} & \textbf{Major implemented methods or intended challenge} \\
\hline
\texttt{sample1.exe} & Emotet simulation & Anti-analysis gating, run-key persistence, encoded configuration recovery, network staging, and process injection; tests whether the workflow can ground persistence and injection claims without overstating full real-world protocol coverage. \\
\texttt{sample2.exe} & QakBot simulation & Mode-driven execution, scheduled-task persistence, resource-backed payload staging, injection, named-pipe coordination, and self-delete support logic; tests careful reporting when resource blobs are intentionally incomplete. \\
\texttt{sample3.exe} & TrickBot simulation & Anti-analysis, integrity checking, registry persistence, encrypted C2 recovery, and staged post-exploitation behavior; tests whether network and banking-style claims stay tied to recovered encrypted configuration and code-level anchors. \\
\texttt{sample4.exe} & WannaCry simulation & Killswitch gating, service persistence, worm-like SMB scanning, ransomware presentation, and resource extraction; tests whether the workflow explains both ransomware and propagation behavior while noting missing-resource limits. \\
\texttt{sample5.exe} & PlugX simulation & Registry persistence, configuration decryption, runtime profile decoding, C2 heartbeat framing, and plugin-style RAT behavior; tests long-horizon configuration recovery and profile interpretation under a more modular RAT-style structure. \\
\end{tabular}
\caption{Blind held-out executable set used in the maintained final evaluation.}
\label{tab:final-heldout-samples}
\end{table}

\begin{table}[p]
\centering
\footnotesize
\renewcommand{\arraystretch}{1.15}
\begin{tabular}{p{0.11\textwidth} p{0.16\textwidth} p{0.25\textwidth} p{0.26\textwidth} p{0.14\textwidth}}
\textbf{Sample} & \textbf{Family simulation} & \textbf{Main obfuscation / evasion} & \textbf{Main malicious behavior cluster} & \textbf{Boundedness to note} \\
\hline
\texttt{sample1.exe} & Emotet & Encoded C2 seeds, RC4-style config decode, FNV-1a API hashing, debugger / VM / analyst-process checks. & Run-key persistence, WinHTTP beaconing, staged module dispatch, and suspended-process injection into \texttt{svchost.exe}. & Network protocol and module loading are simplified rather than full real-world Emotet coverage. \\
\texttt{sample2.exe} & QakBot & XOR-decoded resource-backed payload path, analyst-tool / VM checks, mode-driven branching. & Scheduled-task persistence, resource staging, \texttt{svchost.exe} injection, named-pipe coordination, and self-delete. & The resource payload path is structurally present, but the checked-in repository does not contain the backing blob. \\
\texttt{sample3.exe} & TrickBot & Timing checks, hardware-breakpoint checks, WMI VM detection, NTDLL integrity checks, TLS-callback anti-debugging, PE-header corruption. & Run-key persistence, encrypted C2 recovery, HTTPS check-in, staged module handling, and web-inject-style browser staging. & Module retrieval and browser-facing behavior are simulated rather than a full reflective loader or live webinject stack. \\
\texttt{sample4.exe} & WannaCry & Killswitch gating, service masquerading, packed / stripped final artifact. & Service persistence, backup-deletion commands, ransom-note presentation, target-file enumeration, and SMB-scanning worm behavior. & The source models ransomware and worm flow, but does not implement real file encryption or real EternalBlue exploitation. \\
\texttt{sample5.exe} & PlugX & Encrypted config, XOR runtime-profile decode, staged shellcode-execution path, modular RAT-style command dispatch. & Run-key persistence, TCP heartbeat beaconing, keylogging, screenshot capture, plugin dispatch, and DLL-entry-oriented execution flow. & The shellcode and command handlers are intentionally bounded stubs, and the final artifact is a single EXE rather than a full sideloading package. \\
\end{tabular}
\caption{Final held-out behavior matrix for defense use. This compact summary isolates the main obfuscation/evasion surfaces, the primary malicious behavior cluster, and the most important caveat for each opaque final executable.}
\label{tab:final-heldout-behavior-matrix}
\end{table}

\subsection{Maintained Variant and Configuration Labels}

Unless a row below says otherwise, the baseline refers to the locked configuration used across the thesis experiments: \texttt{pipeline=auto\_triage}, \texttt{architecture=balanced}, \texttt{response\_scope=default}, \texttt{analysis\_hint=default}, \texttt{subagent\_profile=default}, \texttt{worker\_prompt\_shape=default}, \texttt{tool\_profile=full}, and \texttt{model\_profile=repo\_default}. The table below summarizes the maintained variant labels that appear in the archive graphs and final held-out results.

\small
\begin{longtable}{p{0.24\textwidth} p{0.20\textwidth} p{0.38\textwidth} p{0.12\textwidth}}
\caption{Maintained configuration variants referenced throughout the thesis. Each row describes what changes relative to the locked baseline and where that variant appears in the evaluation program.}
\label{tab:evaluation-variants}\\
\textbf{Variant label} & \textbf{Used in} & \textbf{What changes relative to baseline} & \textbf{Dimension} \\
\hline
\endfirsthead
\textbf{Variant label} & \textbf{Used in} & \textbf{What changes relative to baseline} & \textbf{Dimension} \\
\hline
\endhead
\texttt{baseline} & broad, depth, final & No variant override. Serves as the common reference point for all reported deltas and rankings. & reference \\
\texttt{response\_scope:brief} & broad, final & Changes only the response wrapper so the requested final report is more concise and evidence-focused; it does \emph{not} inject new sample-specific hints into the task itself. & prompt framing \\
\texttt{worker\_subagents:single\_generalist} & broad, depth, final & Replaces the default worker-specialist mix with a single generalist worker slot while leaving the rest of the pipeline unchanged. & topology \\
\texttt{worker\_prompt\_shape:artifact\_focused} & broad, depth follow-up, final & Keeps the default worker-role structure but adds an evidence-first worker persona overlay intended to prioritize concrete artifacts and stronger grounding. & prompting \\
\texttt{worker\_prompt\_shape:empty} & broad, depth follow-up, final & Blanks the worker-role prompt while keeping the rest of the workflow fixed, making it a prompt-surface ablation rather than an architecture change. & prompting \\
\texttt{pipeline\_preset:validators\_reporter} & depth follow-up only & Replaces the baseline \texttt{auto\_triage} path with a planner-workers-validators-reporter pipeline so validator review becomes an explicit stage. & pipeline \\
\texttt{architecture\_preset:ghidra\_focused} & final only & Shifts the architecture toward a more Ghidra-centered worker mix and evidence-gathering pattern while preserving the same held-out tasks and tool profile. & architecture \\
\texttt{architecture\_preset:model\_augmented} & final only & Shifts the architecture toward more model-side synthesis and augmentation while holding the benchmark scope, prompt surface, and tool profile constant. & architecture \\
\hline
\end{longtable}
\normalsize

\section{Judge and Prompt Artifacts}

The judge rubric is defined in \path{Testing/config/binary_judge_rubric.json}. It scores six dimensions: \texttt{evidence\_grounding}, \texttt{specificity}, \texttt{technique\_coverage}, \texttt{false\_claim\_control}, \texttt{task\_alignment}, and the down-weighted \texttt{report\_conciseness}. The first five dimensions are each scored from 0 to 5, while conciseness is scored from 0 to 3. Raw totals are normalized to a 0--100 score by dividing by the rubric maximum of \texttt{28}. A run is treated as passing only if the normalized overall score is at least \texttt{70}, \texttt{false\_claim\_control} is at least \texttt{3}, and \texttt{evidence\_grounding} is at least \texttt{3}. This structure is intended to prevent short but weak reports, or confident but unsupported reports, from appearing successful.

The base held-out task prompt is deliberately broad and blind:

\begin{quote}\ttfamily
Analyze this executable as an unknown sample. Recover the main behaviors,
concrete evidence anchors, anti-analysis or packing indicators, and any
persistence, networking, or process-execution behavior while clearly
separating supported findings from uncertainty.
\end{quote}

The response-scope wrappers in \path{Testing/config/response\_scope\_variants.json} do not change the underlying task itself. Instead they wrap that base query with a prefix and suffix that alter how concise or expansive the requested answer should be. For example, the maintained \texttt{brief} variant adds the prefix ``Provide a concise, evidence-grounded reverse-engineering answer'' and the suffix ``Keep the report focused on the most important findings, the strongest supporting evidence, and the main unknowns.'' The analysis-hint wrappers in \path{Testing/config/analysis\_hint\_variants.json} are separate and are currently reserved for future pilot studies. Those wrappers add verify-first investigation leads such as packing, encoded configuration, hashing-based dispatch, or embedded payloads without asserting that such behaviors are present. Table~\ref{tab:evaluation-variants} summarizes which of those maintained wrapper variants are actually exercised in the thesis experiments.

The stage-level prompt surface is split across three files under \path{multi\_agent\_wf/workflow\_config/}: \path{stage\_manager\_prompts.json}, \path{agent\_archetype\_prompts.json}, and \path{stage\_output\_contracts.json}. The first defines the coordinator-level rules for each pipeline stage, the second defines specialist overlays such as \texttt{ghidra\_analyst}, \texttt{control\_flow\_analyst}, or \texttt{string\_analyst}, and the third defines the machine-readable blocks or evidence-bundle contracts that later stages must consume.

\begin{table}[h]
\centering
\small
\begin{tabular}{p{0.16\textwidth} p{0.35\textwidth} p{0.35\textwidth}}
\textbf{Stage} & \textbf{Prompt emphasis} & \textbf{Expected structured output} \\
\hline
Preflight & minimal reconnaissance only; resolve the sample path, basic metadata, and a few worker pivots without deep analysis & short handoff note for planning \\
Deterministic presweeps & bounded reusable collection such as strings, FLOSS, capa, packer checks, UPX attempts, and compact Ghidra metadata & compact reusable sweep bundle \\
Planner & planning only, no fresh analysis; convert the request and shared context into executable work items & JSON work-item block between \texttt{PLANNER\_WORK\_ITEMS\_START} and \texttt{PLANNER\_WORK\_ITEMS\_END} \\
Workers & evidence-backed specialist analysis, role-specific bounded work, optional Ghidra or YARA proposal output & structured worker evidence bundle plus optional machine-readable proposal blocks \\
Validators & review only, no broad new discovery; decide whether the answer is adequate and what should be narrowed or retried & JSON validation decision block between \texttt{VALIDATION\_DECISION\_START} and \texttt{VALIDATION\_DECISION\_END} \\
Reporter & synthesis only; no fresh tool use; final answer must stay tied to validated findings and meaningful control-flow pivots & concise final analyst report and any finalized proposal blocks \\
\end{tabular}
\caption{The main prompt and contract surfaces used by the maintained multi-agent workflow.}
\label{tab:prompt-surfaces}
\end{table}

\section{Configuration and Experiment Provenance}

The workflow-level configuration surface is concentrated in \path{multi\_agent\_wf/workflow\_config/}. The most important files are \path{architecture\_presets.json} for manager and worker topologies, \path{pipeline\_presets.json} for stage sequences, \path{agent\_archetype\_specs.json} and \path{agent\_archetype\_prompts.json} for specialist roles and tool-domain mapping, \path{stage\_manager\_prompts.json} for coordinator instructions, \path{stage\_output\_contracts.json} for machine-readable stage outputs, and \path{worker\_persona\_profiles.json} for worker prompt-shape overlays. The lower-level base prompt scaffolding is defined in \path{base\_prompts.json}.

The evaluation-side configuration surface lives under \path{Testing/config/}. The most important files are \path{presets/*.json} for runnable named presets, \path{experiment\_sweeps.json} for the generic direct-sweep definition, \path{tool\_profiles.json} for MCP availability profiles, \path{response\_scope\_variants.json} and \path{analysis\_hint\_variants.json} for prompt wrappers, \path{binary\_judge\_rubric.json} and \path{prompts/binary\_judge\_prompt.md} for judging, and \path{budget\_guardrails.json} plus \path{model\_cost\_profiles.json} for launch-time cost and safety constraints. Corpus definitions live alongside the sample sets themselves, most importantly \path{Testing/sources/experimental/sample\_manifest.json} and \path{Testing/sources/final\_round/sample\_manifest.json}.

The preset files correspond directly to the main experiment families discussed in this thesis: \texttt{coverage\_broad\_r1\_60usd} for the breadth-oriented pilot, \texttt{sweep\_decoder\_depth\_r3} and \texttt{sweep\_decoder\_depth\_followups\_r3} for the decoder-depth studies, and \texttt{thesis\_final\_eval\_r3\_100usd} for the blind held-out final evaluation. Each preset is self-contained and names its corpus, selected samples or tasks, repetition count, baseline configuration, and sweep families.

When repaired experiments or recovery sessions are cited, the artifact package should preserve three layers of provenance instead of collapsing everything into one synthetic result. First, cite the original experiment directory under \path{Testing/results/experiments/<experiment\_id>/}. Second, cite the recovery session manifest under \path{recovery\_attempts/<session\_id>/recovery\_session.json}, which records which runs and sample-task cases were relaunched. Third, if repaired aggregate figures or summary tables are used, cite the rebuilt repaired experiment directory and its \path{recovery\_rebuild\_manifest.json}. This makes it explicit that a repaired experiment is a derived artifact layered on top of the original sweep, not an in-place replacement of the original raw results. For older pre-refactor development artifacts that still live in the legacy archive, the same principle applies: cite the archived experiment identifier and note when the thesis figure comes from an archive-only repaired or rebuilt result set rather than from the active results tree.

\section{Held-Out Failure Accounting}
\label{app:heldout-failure-accounting}

The main results chapter keeps the unrecovered held-out failures brief so that the repaired final comparisons remain central. This appendix section collects the fuller breakdown for readers who want to inspect the one-shot failure pattern directly. The source artifact for these counts is the original unrecovered final-run summary generated alongside \path{final\_eval\_caveat\_summary.json} and its companion figures.

\begin{table}[h]
\centering
\small
\begin{tabular}{p{0.42\textwidth} p{0.18\textwidth}}
\textbf{Original unrecovered held-out outcome} & \textbf{Count} \\
\hline
Sample-task evaluations attempted & 105 \\
\texttt{completed} & 64 \\
\texttt{worker\_assignment\_failed} & 39 \\
\texttt{analysis\_error} & 2 \\
\end{tabular}
\caption{Task-level outcome totals in the original unrecovered held-out sweep.}
\label{tab:heldout-failure-outcomes}
\end{table}

\begin{table}[h]
\centering
\small
\begin{tabular}{p{0.42\textwidth} p{0.18\textwidth}}
\textbf{Worker-assignment failure class} & \textbf{Count} \\
\hline
Failed worker assignments observed & 58 \\
\texttt{context\_length\_exceeded} & 49 \\
\texttt{invalid\_request\_payload} & 6 \\
\texttt{tool\_retry\_exhausted} & 2 \\
\texttt{usage\_limit\_exceeded} & 1 \\
\end{tabular}
\caption{Assignment-level failure classes in the original unrecovered held-out sweep.}
\label{tab:heldout-failure-causes}
\end{table}

\begin{figure}[p]
\centering
\includegraphics[width=\textwidth]{experiment_results/final_eval_caveat_outcome_by_executable.png}
\caption{Outcome mix by executable in the original unrecovered held-out sweep. The main point of this appendix figure is that the failure burden was not evenly distributed across the five held-out samples.}
\label{fig:appendix-final-eval-outcome-by-executable}
\end{figure}

\begin{figure}[p]
\centering
\includegraphics[width=\textwidth]{experiment_results/final_eval_caveat_worker_failure_causes_by_executable.png}
\caption{Underlying worker-assignment failure causes by executable in the original unrecovered held-out sweep. Most failures came from host-worker context-window overflow on a subset of samples, with a smaller tail of malformed-request, tool-retry, and request-budget failures.}
\label{fig:appendix-final-eval-worker-failure-causes}
\end{figure}

These appendix figures support two narrow conclusions. First, the original one-shot robustness problem was real and was concentrated on a subset of the harder held-out executables rather than evenly distributed across the whole benchmark. Second, the dominant failure mode was infrastructural: most failed worker assignments overflowed the model context window during Ghidra-heavy work. They should not be read as proof that the underlying analytical direction was uniformly poor on those samples.

\section{Detailed Final Per-Executable Panels}
\label{app:final-per-executable-panels}

The main results chapter keeps the heatmap view as the primary executable-level figure because it is the clearest way to show that the repaired final comparison is heterogeneous across samples. The regenerated paged panels below are included for traceability. They preserve the repetition-level score variance for each configuration on each executable without compressing all five samples into one unreadable composite.

\begin{figure}[p]
\centering
\includegraphics[width=\textwidth]{experiment_results/final_eval_per_executable_scores_page_01.png}
\caption{Paged per-executable score panels for the repaired final held-out experiment, page 1. This page covers \texttt{sample1.exe} through \texttt{sample3.exe} and shows the mean score plus repetition-level variance for each maintained configuration.}
\label{fig:appendix-final-per-exec-scores-page1}
\end{figure}

\begin{figure}[p]
\centering
\includegraphics[width=\textwidth]{experiment_results/final_eval_per_executable_scores_page_02.png}
\caption{Paged per-executable score panels for the repaired final held-out experiment, page 2. This page covers \texttt{sample4.exe} and \texttt{sample5.exe} and completes the executable-level detail behind Figure~\ref{fig:final-cross-executable-heatmap}.}
\label{fig:appendix-final-per-exec-scores-page2}
\end{figure}

\section{Reproduction Commands}
\label{app:reproduction-commands}

This section collects the repository commands used to rebuild corpora, run preflight checks, and launch the maintained experiment presets. The main text cites preset names and artifact locations, while the full command lines live here so the body chapters remain focused on methodology and interpretation.

\subsection*{Broad Coverage Pilot}

Rebuild and prepare the experimental corpus:

\begin{verbatim}
make all-exes
python Testing/scripts/prepare_bundles.py --corpus experimental
\end{verbatim}

Preflight and launch the broad pilot:

\begin{verbatim}
python Testing/scripts/run_launch_preset.py \
  --preset coverage_broad_r1_60usd \
  --judge-model openai:gpt-4o-mini \
  --preflight-only

python Testing/scripts/run_launch_preset.py \
  --preset coverage_broad_r1_60usd \
  --judge-model openai:gpt-4o-mini
\end{verbatim}

\subsection*{Focused Decoder-Depth Studies}

The repeated decoder-depth suite:

\begin{verbatim}
python Testing/scripts/run_launch_preset.py \
  --preset sweep_decoder_depth_r3 \
  --judge-model openai:gpt-4o-mini \
  --preflight-only

python Testing/scripts/run_launch_preset.py \
  --preset sweep_decoder_depth_r3 \
  --judge-model openai:gpt-4o-mini
\end{verbatim}

The narrower follow-up depth suite:

\begin{verbatim}
python Testing/scripts/run_launch_preset.py \
  --preset sweep_decoder_depth_followups_r3 \
  --judge-model openai:gpt-4o-mini \
  --preflight-only

python Testing/scripts/run_launch_preset.py \
  --preset sweep_decoder_depth_followups_r3 \
  --judge-model openai:gpt-4o-mini
\end{verbatim}

\subsection*{Final Held-Out Evaluation}

Rebuild and prepare the final corpus:

\begin{verbatim}
make deepseek-final
python Testing/scripts/prepare_bundles.py --corpus final_round
python Testing/scripts/run_launch_doctor.py --corpus final_round --bundle-only
\end{verbatim}

Preflight and launch the held-out evaluation:

\begin{verbatim}
python Testing/scripts/run_launch_preset.py \
  --preset thesis_final_eval_r3_100usd \
  --judge-model openai:gpt-4o-mini \
  --preflight-only

python Testing/scripts/run_launch_preset.py \
  --preset thesis_final_eval_r3_100usd \
  --judge-model openai:gpt-4o-mini \
  --live-view \
  --max-concurrent-child-runs 9 \
  --max-concurrent-repetitions 3 \
  --enable-budget-guardrails
\end{verbatim}

Preflight-only artifacts are written under \path{Testing/results/preflight/}. Executed experiments write canonical run catalogs, child runs, summaries, and figures under \path{Testing/results/experiments/<experiment\_id>/}.
